\hypertarget{conditions-settlement-choices-and-the-earnings-that-each-share-receives}{%
\section{Conditions, settlement choices and the earnings that each share receives}\label{conditions-settlement-choices-and-the-earnings-that-each-share-receives}}

The following conditional methods use IFRS 2 and IAS 33 in the retained September 2023 EU compilation. They are separate from the Apple contractual examples and do not determine their US GAAP treatment. Employee counts, fair values, service patterns and participation rights below are explicit synthetic inputs. Valuing the awards and admitting an entity's applicable rules remain distinct requirements.

\hypertarget{a-failed-condition-does-not-always-reverse-an-expense}{%
\subsection{A failed condition does not always reverse an expense}\label{a-failed-condition-does-not-always-reverse-an-expense}}

An employee's failure to complete required service differs from a share-price target that is not reached. For employee equity awards, IFRS~2.10--21 \citep[IFRS~2.10--23]{eu20231803} measures the services using grant-date fair value. Service and non-market performance conditions affect the estimated number of instruments expected to vest; market conditions enter fair value. Revise the former estimate as information changes and use the actual eligible number at vesting. If a market condition fails but the other vesting conditions are met, the service expense remains. Freezing the quantity at grant and subsequently marking every equity award to the current share price would get both sides of this distinction wrong.

The award record therefore separates grant and measurement dates, recipient and employing entity, each service/performance/non-vesting condition and who can choose to fulfil it, service already received, expected eligible quantity, grant-date valuation and its incorporated conditions, settlement rights, cumulative expense, equity and any liability. A contract's use of the word performance does not decide whether the condition is a market condition. Missing terms or fair value remain unresolved rather than selecting the most convenient expense pattern. Services are expensed in these examples; qualifying goods or services that create an asset require their own recognition treatment.

Give 100 employees ten equity instruments each, with grant-date fair value 3 per instrument and three years of equal service. At the first year-end, 90 employees are expected to complete service; at the second, 80; at vesting, 75 have completed it. The cumulative expenses are

\[90(10)(3)/3=900,\qquad 80(10)(3)(2/3)=1600,\qquad75(10)(3)=2250.\]

The annual charges are therefore 900, 700 and 650, credited to award equity. If an additional non-market performance condition ultimately fails for all awards, cumulative expense becomes zero and previously recognised expense reverses. In a separate award whose supplied grant-date fair value of 3 already incorporates a market hurdle, failure of that hurdle does not reverse the 2250 for the 75 employees who satisfy the other conditions. Equal supplied fair values in these examples do not imply that different contracts have the same economic value. Once an equity award has vested, later non-exercise does not reverse total equity under paragraph 23; a transfer within equity remains possible.

Cancellation also differs from forfeiture. Paragraphs 27--28 \citep[IFRS~2.27--28A]{eu20231803} preserve the minimum original grant-date expense, subject to failure of the specified non-market vesting conditions, and accelerate unrecognised service expense on cancellation during vesting. For one grant with total grant-date value 300, one of three equal service years gives expense 100. An entity cancellation that is not a forfeiture accelerates the remaining 200. If it pays 330 when the repurchased award's fair value is 320, debit equity 320 and expense 10, and credit cash 330. Total service/cancellation expense is 310. Starting from cash and capital 1000 leaves cash and equity 670 after all entries. Calling the event a service forfeiture would wrongly reverse the first 100.

Non-vesting conditions enter fair value under paragraph 21A. However, when the entity or counterparty can choose whether to meet such a condition, failure to do so during vesting is a cancellation under paragraph 28A. Consequently the condition's classification and the party's control of it must be preserved together. Beneficial modifications require their incremental value analysis under paragraph 27 and Appendix B; this cancellation example does not replace that calculation.

\hypertarget{cash-obligations-change-as-the-awards-value-changes}{%
\subsection{Cash obligations change as the award's value changes}\label{cash-obligations-change-as-the-awards-value-changes}}

For a cash-settled award, IFRS~2.30--33D \citep[IFRS~2.30--33D]{eu20231803} requires remeasurement at each reporting date and settlement. Service and non-market vesting conditions affect quantity; market and non-vesting conditions enter the current valuation. Apply the service fraction to the current eligible quantity and fair value, then charge the difference from the previously recognised liability. The final cumulative expense equals cash paid, unlike the fixed grant-date measurement of an equity award.

Use the preceding employee quantities and ten rights per employee, but admit cash-settled fair values of 3, 6 and 8 at the three year-ends. Cumulative liabilities are 900, 3200 and 6000. The annual expenses are 900, 2300 and 2800. After vesting, a settlement-date value of 9 raises the liability to 6750 and adds expense 750. Paying 6750 removes the liability. A model that freezes the original value of 3 would leave expense 2250 and an unrecorded cash obligation of 4500. The changing values are supplied, not estimated from this example.

\hypertarget{who-chooses-settlement-matters}{%
\subsection{Who chooses settlement matters}\label{who-chooses-settlement-matters}}

If the employee can demand cash or shares, paragraphs 35--40 \citep[IFRS~2.35--43]{eu20231803} treat the award as a compound instrument. Measure the debt component first, then the equity component taking account of the cash right that must be surrendered for shares. Do not add the standalone values of both alternatives. Over the service period, remeasure the debt component under the cash-settled method and recognise the grant-date equity component under the equity-settled method.

For an immediately vested award, admit debt value 100 and incremental equity value 20. Recognise expense 120, liability 100 and equity 20. At settlement, remeasure the liability to 130 and charge expense 30. If the employee takes cash, pay 130 and retain the original 20 within equity. If the employee takes shares, transfer the liability 130 directly to equity, leaving total award equity 150. Both branches have cumulative expense 150; only the cash branch has a cash outflow. The values are an admitted compound valuation, not a claim that standalone alternative values uniquely determine the equity part.

If the entity chooses, paragraphs 41--43 first ask whether it nevertheless has a present obligation to settle in cash. Commercially insubstantial share settlement, a past practice or stated policy of paying cash, or normally paying cash whenever requested can create that obligation. A written share option therefore does not by itself establish an equity award. If no cash obligation exists, use equity-settled accounting. At settlement, cash is an equity repurchase, except that choosing the higher-valued alternative creates expense for the excess. For an immediately vested grant with admitted value 100, settlement cash 130 and alternative shares worth 120 require a debit to equity of 120, additional expense 10 and cash credit 130. A 130 expense would count the repurchase as remuneration again.

These cash choices are distinct from statutory withholding. Under paragraphs 33E--33H, an otherwise equity award can retain equity classification when shares are withheld to satisfy a legally imposed employee-tax withholding obligation. Preserve the legal obligation and required quantity. The exception does not cover an arrangement without that obligation or excess shares paid out in cash. It cannot justify classifying every net-settled award as equity.

\subsection{A beneficial amendment adds a second service-cost clock}
\label{sec:award-modifications}

Cancellation accelerates an award's remaining cost. A beneficial amendment
instead normally preserves the original grant-date cost and adds value measured
at the modification date. IFRS~2.B42--B44 requires the original minimum, subject
to failure of the relevant original non-market vesting conditions, and recognizes
the beneficial change. The comparison uses the award's value immediately before
and after modification at the same date
\citep[IFRS~2.27, B42--B44]{eu20231803}. A higher current value is not a
replacement grant-date value.

Grant 100 instruments with fair value 3 each for three years of equal service.
All ultimately vest. After the first year, a beneficial repricing raises their
supplied fair value from 2 to 4 per instrument, with the original vesting date
unchanged. Original total cost is $100(3)=300$. Incremental cost is
$100(4-2)=200$, spread over the two remaining service years. The first year's
expense is 100; each later year records original cost 100 plus incremental
cost 100. Total expense and award equity are 500. Replacing the grant with its
current value of 400 would omit 100; subtracting the grant-date value 3 from
the modification-date value 4 would also use the wrong comparison. Each annual
entry debits service expense and credits award equity. With no cash or other
movements, the increase in that equity component offsets the reduction in
retained earnings, leaving total equity unchanged.

The two service periods need not end together. In a separate modification,
leave the original 100 instruments unchanged and add 20 instruments worth 4
each that require only the second year's service. IFRS~2.B43(b) measures the
additional 80 at the modification date over their own service period. Annual
charges are now 100, 180 and 100, totaling 380. Spreading the additional 80
over both remaining original years would understate year 2 by 40. This example
does not shorten the original award's service condition: a beneficial change
to that condition must itself be reflected under B43(c). An adverse value-only
change generally leaves the original 300 cost, while reducing the number of
instruments invokes the cancellation rule for the affected portion
\citep[IFRS~2.B43(b)--(c), B44]{eu20231803}.

A change from cash settlement to equity settlement follows a distinct rule.
Suppose two of three equal service years have passed. The cash award liability,
remeasured at the modification date, is 180. Replacement equity has a total
modification-date fair value of 300 and retains the third year's service
requirement. IFRS~2.B44A recognizes $300(2/3)=200$ in equity for services
already received, removes the liability 180 and charges the difference 20
immediately. Debit liability 180 and expense 20; credit award equity 200.
The remaining year charges 100 to expense and equity. Prior cash-award expense
180 plus the modification charge 20 plus later service 100 gives lifetime
expense 300, with no cash payment and no remaining liability. If the correctly
remeasured liability had instead been 220, the immediate difference would be
an expense credit of 20. It is not automatically an extra positive cost
\citep[IFRS~2.B44A--B44C]{eu20231803}.

Keep the original grant, every modification's before-and-after valuation,
eligible quantity, original and modified service conditions, settlement
classification, cumulative expense and component balances together. Apply
each event once and recognize the change from the previously recorded
cumulative amounts. B44B requires the modified vesting period when applying
the cash-to-equity rule; B44C extends it to an identified replacement for a
cancelled cash award. Supplied fair values in these examples do not provide
an option-pricing method. Missing valuation or legal terms prevent application
without erasing the original award or assuming that it costs zero.


\hypertarget{participation-allocates-the-numerator-before-division}{%
\subsection{Participation allocates the numerator before division}\label{participation-allocates-the-numerator-before-division}}

IAS~33.A13--A14 \citep[IAS~33.A13--A16]{eu20231803} allocates earnings according to the actual dividend and undistributed-earnings rights of the classes. First deduct the declared and contractually required dividends from the period's allocable result. Allocate the remainder as if it were distributed, using the participation formula and any caps. Add each class's dividends back to its allocation, then divide by the instruments to which that allocation belongs. A class's right to profits does not establish a right or obligation to share losses; preserve the actual terms. Unknown participation rights block the allocation.

Suppose parent-attributable profit is 120, 100 ordinary shares participate with weight one, and ten nonconvertible participating shares have weight two. Assume all are outstanding throughout the period and there are no senior claims, caps or other instruments. Declared dividends are 20 and 4. Remaining profit 96 allocates over total weight \(100+2(10)=120\): ordinary shares receive 80 and participating shares receive 16. Adding dividends gives class earnings 100 and 20, so per-share earnings are 1 and 2. Allocating all 120 to the ordinary shares would overstate their EPS as 1.2. Dividends are an allocation of this profit, not a second expense. A capped or senior participation formula needs its own piecewise allocation; the simple weights cannot replace it.

Diluted EPS then applies the relevant potential-share rules and dilution tests; it is not the same calculation as award expense. Under paragraphs 47A--48 \citep[IAS~33.47A--60]{eu20231803}, future goods or services form part of the assumed proceeds for employee options. A separate case with profit 100, 100 ordinary shares, ten options, cash exercise price 5, period-average future service value 2 per option and average share price 10 has proceeds 70. Seven shares are assumed repurchased, so incremental shares are three and diluted EPS is \(100/103\). Omitting future services would instead add five. This comparison assumes the awards are outstanding for the whole period and pass the applicable dilution test; performance-conditioned awards also require paragraphs 52--57's period-end contingency assessment.

Under paragraphs 58--60, an entity cash/share choice presumes share settlement if dilutive, whereas a holder choice uses the more dilutive outcome. Required numerator adjustments preserve the profit effect of assumed equity treatment. For a separate nonparticipating contract, let reported profit be 90 on 100 shares after liability-remeasurement expense 10 that would be absent under wholly equity classification, and let assumed share settlement add 20 shares. With no tax effect or other changes, cash gives EPS 0.9 and share settlement gives \(100/120=5/6\). The share outcome is more dilutive. Adding shares while leaving profit at 90 would produce 0.75 and compute the wrong assumed world. Continuing-operations control, instrument ordering, time weights and antidilution remain the earlier EPS method's requirements, not assumptions waived by this example.
